\section{Признаки Дирихле и Абеля}\label{sec-series-03}

В обоих признаках \(u_n\) могут быть комплекснозначными, а \(v_n\) --- \textbf{вещественными}, поскольку требуется монотонность по \(n\) в каждой точке \(x\).

\begin{theorem}[Признак Дирихле]\label{thm-dirichlet}

Пусть:

\begin{enumerate}
\def\labelenumi{\arabic{enumi}.}
\tightlist
\item
  Частичные суммы \(U_N(x)=\sum_{k=1}^N u_k(x)\) равномерно ограничены: \(|U_N(x)|\le C\) при всех \(N\) и \(x\in\Omega\).
\item
  Для каждого \(x\) последовательность \(v_n(x)\) монотонна по \(n\).
\item
  \(v_n\rightrightarrows0\) на \(\Omega\).
\end{enumerate}

Тогда \(\sum u_n(x)v_n(x)\) равномерно сходится на \(\Omega\).

\end{theorem}

\begin{proof}

Для хвоста от \(n+1\) до \(m\) обозначим \(A_k=\sum_{j=n+1}^k u_j\). Тогда \(|A_k|\le2C\). Преобразование Абеля: \[
\sum_{k=n+1}^m u_kv_k=A_mv_m+\sum_{k=n+1}^{m-1}A_k(v_k-v_{k+1}).
\] Монотонная последовательность, сходящаяся к нулю, имеет постоянный знак и убывающий модуль. Поэтому \[
\left|\sum_{k=n+1}^m u_k(x)v_k(x)\right|
\le2C\left(|v_m(x)|+\sum_{k=n+1}^{m-1}|v_k(x)-v_{k+1}(x)|\right)
=2C|v_{n+1}(x)|.
\] Правая часть стремится к нулю равномерно по \(x\), независимо от \(m>n\). Применяем теорему \ref{thm-series-cauchy}.

\end{proof}

\begin{theorem}[Признак Абеля]\label{thm-abel-test}

Пусть \(\sum u_n(x)\) равномерно сходится на \(\Omega\), \(v_n(x)\) монотонна по \(n\) для каждого \(x\), а \(|v_n(x)|\le C\) для всех \(n\) и \(x\). Тогда \(\sum u_n(x)v_n(x)\) равномерно сходится на \(\Omega\).

\end{theorem}

\begin{proof}

По критерию Коши при достаточно большом \(n\) все хвостовые суммы \(A_k=\sum_{j=n+1}^k u_j\) удовлетворяют \(|A_k(x)|<\eta\) одновременно для всех \(k>n\) и \(x\). То же преобразование Абеля и монотонность дают \[
\left|\sum_{k=n+1}^m u_kv_k\right|
\le\eta\left(|v_m|+|v_{n+1}-v_m|\right)\le3C\eta.
\] Выбираем \(\eta=\varepsilon/(3C)\) при \(C>0\); случай \(C=0\) тривиален.

\end{proof}
